\chapter{Randomized Least Squares} \label{ch:randls}

This chapter develops the sketching-based least-squares methods that form the algorithmic core of the thesis. The starting point is the overdetermined problem
\[
    \min_{x \in \R^d} \|Ax-b\|_2,
\]
where \(A \in \R^{n \times d}\) has many more rows than columns. Interpreted geometrically, one seeks the point in the range of the linear map represented by \(A\) that is closest to \(b\). When \(n\) is large, the dominant cost lies not in the abstract existence theory, but in the repeated dense linear algebra required to solve or approximate the system.

The randomized approach replaces the original problem by a smaller one. Given a sketching matrix \(S \in \R^{s \times n}\) with \(s \ll n\), one solves
\[
    \min_{x \in \R^d} \|SAx-Sb\|_2.
\]
If \(S\) approximately preserves the geometry of the range of \(A\), then the reduced problem is expected to return a solution whose residual is close to optimal while being cheaper to compute. In operator terms, \(S\) should act almost isometrically on the subspace that carries the least-squares information.

\section{Gaussian Sketching Model}

The first sketch family considered in this thesis is the Gaussian sketch. Its entries are independent normal random variables with variance \(1/s\), so that
\[
    S_{ij} \sim \mathcal{N}(0, 1/s).
\]
This scaling is chosen so that, in expectation, the sketch acts approximately like an isometry on Euclidean norms. In particular, for a fixed vector \(x \in \R^n\), one has the heuristic identity
\[
    \mathbb{E}\|Sx\|_2^2 = \|x\|_2^2,
\]
which explains why Gaussian sketches are a natural baseline in both theory and computation. The relevant point is not merely distributional symmetry, but the fact that the sketch preserves the metric structure of vectors in expectation and, with high probability, approximately preserves it on low-dimensional subspaces.

The Gaussian model is mathematically convenient because it avoids combinatorial structure and makes concentration arguments comparatively clean. For that reason, it is also the first model implemented in the numerical experiments and in the accompanying code.

\section{A First Experiment}

The initial computational study in this thesis is deliberately simple. A random matrix \(A \in \R^{n \times d}\) is generated together with a reference solution \(x_\star\), and the right-hand side is taken as either
\[
    b = Ax_\star
\]
in the consistent case or
\[
    b = Ax_\star + e
\]
in the noisy case, where \(e\) is a perturbation vector. For each sketch size \(s\), several independent Gaussian sketches are drawn, the sketched least-squares problems are solved, and the following quantities are recorded:
\begin{itemize}
    \item the mean absolute error of the averaged solution,
    \item the residual ratio relative to the unsketched least-squares solution,
    \item the relative solution error with respect to a reference solve,
    \item the condition number of the sketched matrix \(SA\).
\end{itemize}

This experiment plays an important organizational role. Mathematically, it gives a concrete setting in which subspace embedding heuristics can be discussed. Computationally, it provides the first benchmark that is shared across the Octave prototype and the R package implementation. In that sense, it is the smallest example in the repository where the operator-theoretic picture, the numerical conditioning, and the measured computation are meant to line up exactly.

\section{Why This Experiment Matters}

Even before formal guarantees are introduced, the experiment highlights the main questions that drive the later chapters. How large must the sketch dimension \(s\) be before the residual becomes reliably close to optimal? How does the conditioning of the compressed operator \(SA\) change as \(s\) varies? Which error metrics stabilize quickly, and which remain sensitive? Once mixed precision is introduced, the same framework also makes it possible to ask which arithmetic stages tolerate lower precision and which do not.

The main topics to be developed further are Gaussian sketches, structured sketches such as SRHT, residual and solution guarantees, and the conditioning of the reduced systems.
